\begin{figure}[tb]
\centering
\begin{tikzpicture}[
  x=1cm,y=1cm,
  terminal/.style={circle,fill=black,inner sep=2.3pt},
  contracted/.style={circle,draw=black,fill=white,line width=.7pt,inner sep=3pt},
  block/.style={draw=black!55,rounded corners=3pt},
  cutshade/.style={fill=blue!12,draw=none},
  edge/.style={draw=black!50,line width=.55pt},
  every node/.style={font=\small}
]
\path[cutshade] (0.25,1.3)--(2.55,1.3)--(2.55,.2)--(4.55,.2)
               --(4.55,2.25)--(.25,2.25)--cycle;
\draw[block] (.15,.1) rectangle (2.25,2.35);
\draw[block] (2.45,.1) rectangle (4.65,2.35);
\node at (1.2,2.68) {$A$};
\node at (3.55,2.68) {$\overline A$};
\node[terminal] (a1) at (1.2,1.8) {};
\node[terminal] (a2) at (1.2,.65) {};
\node[terminal] (b1) at (3.55,1.8) {};
\node[terminal] (b2) at (3.55,.65) {};
\draw[edge] (a1)--(a2);
\draw[edge] (a1)--(b1);
\draw[edge] (a1)--(b2);
\draw[edge] (a2)--(b1);
\draw[edge] (a2)--(b2);
\draw[edge] (b1)--(b2);
\node at (2.4,-.3) {lifted side: three terminals};

\draw[-{Stealth[length=2mm]},line width=.7pt]
  (4.95,1.2)--(6.35,1.2);
\node[align=center,font=\footnotesize] at (5.65,1.8)
  {contract\\$\overline A$};

\path[cutshade] (6.7,1.3) rectangle (10.8,2.25);
\draw[block] (6.6,.1) rectangle (8.7,2.35);
\node at (7.65,2.68) {$A$};
\node[terminal] (c1) at (7.65,1.8) {};
\node[terminal] (c2) at (7.65,.65) {};
\node[contracted] (z) at (10.2,1.8) {};
\draw[edge] (c1)--(c2);
\draw[edge] (c1)--(z);
\draw[edge] (c2)--(z);
\node[align=center,font=\footnotesize] at (10.2,.7)
  {nonterminal\\even number of\\original terminals};
\node at (8.7,-.3) {child side: one terminal};
\end{tikzpicture}
\caption{Lifting through an even-terminal contraction.
Filled dots are terminals and the open dot represents the contracted
block.  In this schematic instance each side initially has two terminals.
The shaded child side lifts to the shaded parent side by taking the
entire contracted block, preserving odd parity.  Parallel crossing
capacities are added, so the cut capacity is preserved as well.}
\label{d:fig-odd-cut-lift}
\end{figure}
